\documentclass[a4paper,12pt]{article}

\usepackage[T2A]{fontenc}
\usepackage[utf8]{inputenc}
\usepackage[russian]{babel}

\usepackage{amsmath, amssymb}
\usepackage{geometry}
\usepackage{fancyhdr}
\usepackage{enumitem}

\geometry{margin=2cm}
\setlength{\parskip}{6pt}
\setlength{\parindent}{0pt}

\pagestyle{fancy}
\fancyhf{}
\cfoot{\thepage}
\rhead{8 класс}

\begin{document}

\begin{center}
    {\Large \textbf{Делимость. Теорема Вильсона, Эйлера, Ферма.}}\\[0.5em]
\end{center}
 
\begin{enumerate}[label=\textbf{Задача \arabic*.}]
\item Решите систему сравнений $\begin{cases} 6x + 5y \equiv 1 \pmod{11}, \\ 4x + 3y \equiv 2 \pmod{11}. \end{cases}$
\item Существует ли степень тройки, заканчивающаяся на 0001? 
\item Докажите, что если p – простое число и  $1 \leq k \leq p - 1$,  то $\binom{n}{k}$ делится на p ($\binom{n}{k}$ - число сочетаний из n по k). 
\item Докажите, что  $5^{120} - 1$  делится на 143.
\item Докажите, что $(p - 2)! \equiv 1 \pmod{p}$ при любом простом $p$.
\item С помощью теоремы Вильсона докажите, что среди чисел вида $n! + 1$ бесконечно много составных.
\item Докажите, что среди любых n последовательных натуральных чисел есть число, которое делится нацело на n.
\item Докажите, что число $\frac{10^{2026} - 1}{9}$  делится нацело на 2027.
\item Докажите, что число является квадратом целого числа тогда и только тогда, когда оно имеет нечетное число натуральных делителей.
\item Натуральное число n не имеет собственных делителей, больших 1 и не превосходящих $\sqrt{n}$. Докажите, что n - простое число.
\item Пусть $\phi(n)$ — количество натуральных чисел от 1 до n, взаимно простых с n. Докажите, что $\phi(n)$ - четное число при $n > 2$.
\item Докажите, что существует сколь угодно много подряд идущих составных чисел (полезно рассмотреть факториал числа).
\end{enumerate}
\end{document}